\documentclass[11pt]{article}
\usepackage[spanish]{babel}
\usepackage[T1]{fontenc}
\usepackage[utf8]{inputenc}
\usepackage{lmodern}
\usepackage[margin=0.8in]{geometry}
\usepackage[table]{xcolor}
\usepackage{array}
\usepackage{tabularx}
\usepackage{ragged2e}
\usepackage{booktabs}
\usepackage{pdflscape}
\usepackage{enumitem}
\usepackage{graphicx}

\setlength{\parindent}{0pt}
\setlength{\parskip}{0.45em}
\setlength{\tabcolsep}{4pt}
\renewcommand{\arraystretch}{1.18}

\definecolor{headbg}{RGB}{220,230,241}
\definecolor{subbg}{RGB}{242,242,242}

\newcolumntype{Y}{>{\RaggedRight\arraybackslash}X}

\begin{document}

\begin{center}
{\Large\textbf{Matriz de Consistencia}}\\[0.4em]
{\normalsize\textbf{Grupo 6}}\\[0.3em]
{\normalsize Modelo de aprendizaje profundo con mecanismos de atención para mejorar la calidad del servicio de guía virtual en servicios de guía virtual con aplicaciones móviles e IoT con señales Wi-Fi/BLE en museos y espacios culturales, 2026}
\end{center}

\textbf{Línea de investigación:} Inteligencia Artificial y Comunicación Inalámbrica.

\textbf{Problema general:} ¿Cómo mejorar la calidad del servicio de guía virtual en servicios de guía virtual con aplicaciones móviles e IoT con señales Wi-Fi/BLE en museos y espacios culturales, 2026?

\textbf{Objeto de estudio:} Servicios de guía virtual.

\textbf{Objetivo general:} Mejorar la calidad del servicio de guía virtual en servicios de guía virtual con aplicaciones móviles e IoT con señales Wi-Fi/BLE en museos y espacios culturales, 2026.

\textbf{Hipótesis general:} La aplicación de un modelo de aprendizaje profundo con mecanismos de atención mejora significativamente la calidad del servicio de guía virtual en servicios de guía virtual con aplicaciones móviles e IoT con señales Wi-Fi/BLE en museos y espacios culturales, 2026.

\begin{landscape}
\footnotesize

\rowcolors{2}{white}{subbg}
\resizebox{\linewidth}{!}{%
\begin{tabular}{|p{5.4cm}|p{5.4cm}|p{5.2cm}|p{7.0cm}|}
\hline
\rowcolor{headbg}
\textbf{Problema específico 1} & \textbf{Objetivo específico 1} & \textbf{Hipótesis específica 1} & \textbf{Variables, dimensiones e indicadores relacionados} \\
\hline
¿De qué manera la variabilidad del RSSI, el multipath y la heterogeneidad de dispositivos afectan la precisión de localización indoor en servicios de guía virtual?
&
Analizar el efecto de la variabilidad del RSSI, el multipath y la heterogeneidad de dispositivos sobre la precisión de localización indoor.
&
La variabilidad del RSSI, el multipath y la heterogeneidad de dispositivos reducen significativamente la precisión de los sistemas de localización indoor basados en Wi-Fi/BLE.
&
\textbf{Variable dependiente:} Calidad del servicio de guía virtual.\\
\textbf{Dimensiones:}
\begin{itemize}[leftmargin=1.1em, itemsep=0.15em, topsep=0.2em]
\item precisión de localización;
\item oportunidad de activación;
\item pertinencia contextual;
\item continuidad del servicio.
\end{itemize}
\textbf{Indicadores:}
\begin{itemize}[leftmargin=1.1em, itemsep=0.15em, topsep=0.2em]
\item porcentaje de activaciones correctas;
\item tiempo de respuesta de la guía;
\item tasa de error por sala o zona;
\item porcentaje de sesiones con guía contextual correcta.
\end{itemize}
\\
\hline
\end{tabular}%
}

\vspace{0.8em}

\resizebox{\linewidth}{!}{%
\begin{tabular}{|p{5.4cm}|p{5.4cm}|p{5.2cm}|p{7.0cm}|}
\hline
\rowcolor{headbg}
\textbf{Problema específico 2} & \textbf{Objetivo específico 2} & \textbf{Hipótesis específica 2} & \textbf{Variables, dimensiones e indicadores relacionados} \\
\hline
¿En qué medida un modelo de aprendizaje profundo con mecanismos de atención puede mejorar la calidad del servicio de guía virtual en servicios de guía virtual con aplicaciones móviles e IoT con señales Wi-Fi/BLE?
&
Diseñar e implementar un modelo de aprendizaje profundo con mecanismos de atención para mejorar la calidad del servicio de guía virtual en servicios de guía virtual con aplicaciones móviles e IoT con señales Wi-Fi/BLE.
&
Un modelo de aprendizaje profundo con mecanismos de atención mejora la calidad del servicio de guía virtual frente a enfoques base en servicios de guía virtual con aplicaciones móviles e IoT con señales Wi-Fi/BLE.
&
\textbf{Variable independiente:} Modelo de aprendizaje profundo con mecanismos de atención.\\
\textbf{Dimensiones:}
\begin{itemize}[leftmargin=1.1em, itemsep=0.15em, topsep=0.2em]
\item arquitectura del modelo;
\item mecanismo de atención;
\item robustez frente a ruido y heterogeneidad;
\item capacidad de generalización.
\end{itemize}
\textbf{Indicadores:}
\begin{itemize}[leftmargin=1.1em, itemsep=0.15em, topsep=0.2em]
\item accuracy de clasificación por zona o sala;
\item error medio de localización;
\item F1-score;
\item estabilidad entre dispositivos o escenarios.
\end{itemize}
\\
\hline
\end{tabular}%
}

\vspace{0.8em}

\resizebox{\linewidth}{!}{%
\begin{tabular}{|p{5.4cm}|p{5.4cm}|p{5.2cm}|p{7.0cm}|}
\hline
\rowcolor{headbg}
\textbf{Problema específico 3} & \textbf{Objetivo específico 3} & \textbf{Hipótesis específica 3} & \textbf{Alcances y relación con el estudio} \\
\hline
¿Cómo impacta la mejora de la precisión de localización indoor en la calidad del servicio de guía virtual en museos y espacios culturales?
&
Evaluar el impacto de la mejora de la precisión de localización indoor sobre la calidad del servicio de guía virtual.
&
Una mayor precisión de localización indoor mejora la activación contextual, la oportunidad y la pertinencia del servicio de guía virtual.
&
\textbf{Objeto de estudio:} Servicios de guía virtual.\\
\textbf{Alcance espacial:} Museos y espacios culturales.\\
\textbf{Alcance temporal:} 2026.
\\
\hline
\end{tabular}%
}

\vspace{0.9em}

\resizebox{\linewidth}{!}{%
\begin{tabular}{|p{5.2cm}|p{18.2cm}|}
\hline
\rowcolor{headbg}
\textbf{Diseño metodológico preliminar} & \textbf{Descripción} \\
\hline
Tipo de investigación & Aplicada. \\
\hline
Nivel & Explicativo y experimental. \\
\hline
Método & Cuantitativo. \\
\hline
Diseño & Cuasi experimental con comparación entre baseline y modelo propuesto. \\
\hline
Población & Registros de señales Wi-Fi/BLE y sesiones de guía virtual en espacios indoor. \\
\hline
Muestra & Subconjunto de huellas o dataset de entrenamiento, validación y prueba. \\
\hline
Técnicas de recolección & Captura de fingerprints, registro de RSSI y evaluación del servicio. \\
\hline
Técnicas de análisis & Métricas de clasificación/localización y comparación estadística descriptiva. \\
\hline
\end{tabular}%
}

\end{landscape}

\end{document}
